\section{Related Work}\label{related-work}

\subsection{RAG Retrieval Evaluation and Embedding Benchmarking}\label{rag-retrieval-evaluation-and-embedding-benchmarking}

RAG systems depend on retrieval quality to maintain answer grounding and downstream reliability, so retrieval design remains a central determinant of end-task behavior. Foundational work establishes this dependency, and later survey evidence frames retrieval as a multi-component design space spanning query processing, retrieval, reranking, and generation-time integration. \citep{Lewis_RAG_2021, Zhu_LLMIR_Survey_2025}

Recent applied studies reinforce this point in domain-specific settings. Knowledge-graph-augmented RAG pipelines for cybersecurity and customs support report that retrieval-side design choices materially affect answer quality and contextual relevance. \citep{Gao_KG_RAG_Power_2025, Hu_ICCARAG_2025} QuIM-RAG further demonstrates that retrieval-oriented query matching strategies can improve question-answering effectiveness relative to simpler baseline retrieval flows. \citep{Saha_QuIMRAG_2024}

Related evidence from question-retrieval research also indicates that metadata-aware retrieval strategies can improve ranking quality in practical question-answering environments, which is consistent with the broader claim that retrieval quality is highly configuration-sensitive. \citep{Ghasemi_Metadata_QR_2024}

Taken together, these studies motivate retrieval-aware system design, but they are typically optimized for one task setting or one domain-specific application pipeline. In contrast, this paper focuses on controlled cross-dataset embedding-model selection under a fixed evaluation protocol, with reusable benchmark artifacts and reporting outputs intended for practitioner-facing local deployment decisions.

\subsection{Preliminary Conference Study (Short Context)}\label{preliminary-conference-study-short-context}

Our earlier conference paper (\texttt{publishedConferencePaper.pdf}) introduced a local automated benchmark pipeline and demonstrated it on a single technical English corpus. That preliminary work established feasibility for end-to-end automated evaluation with synthetic queries, LLM-based qrels, and multi-model retrieval metrics. \citep{Karakaya_Conference_2025}

\subsection{Extension in the Current Journal Work}\label{extension-in-the-current-journal-work}

This journal paper extends the preliminary work in scope and practical selection value: it uses a frozen cross-dataset setup with \texttt{education\_eng}, \texttt{education\_tr}, and \texttt{law\_tr}, and evaluates 5 embedding models under one consistent pipeline configuration. The dataset sources are explicitly documented (TED regulations/guidelines for education datasets and BDDK regulations for law dataset, retrieved on 01/01/2026). The extension adds aggregate cross-dataset ranking, consistency analysis, and quality-speed analysis, plus automatic benchmarking reports across various metrics, to support cross-dataset comparison for dataset-specific model selection. \citep{Aggregate_Report_2026_03_05, Raw_Education_Eng_2026_03_05, Raw_Education_Tr_2026_03_05, Raw_Law_Tr_2026_03_05, TED_Regulations, BDDK_Mevzuat}

In short, the conference paper is context only; the journal evidence is driven by the current multi-dataset artifacts and emphasizes that the pipeline is reusable by practitioners for choosing the most suitable embedding model for application requirements. \citep{Karakaya_Conference_2025, GitHub_RagRetrievalBench}
